\section{第一轮独立评估协议与统计解释}
\label{sec:experimental-protocol}

\subsection{数据用途分离与冻结边界}

实验区分方法发现、开发验证、最终点估计比较以及预测区域校准。研究工具箱只提供发现集与验证集接口：前者用于提出和修订假设，后者允许有限次评价，因此也属于开发过程，不能称为最终无偏测试。方法冻结后，\texttt{benchmark\_v2\_1.py}读取冻结文件及登记表，写入协议、方法配置、源代码摘要和运行环境，然后生成独立评估案例。表\ref{tab:dataset-protocol}列出当前默认协议；运行时参数如有改变，应以该次保存的\texttt{protocol.json}为准。

评估器在研究尚未冻结、任何正式测试案例尚未生成时接受一次预先代码审查修订：当最终规格为旧基线简称时，不生成新组合方法的非法组件消融；正常组合方法分支完全不变。正式首次测试统一采用\texttt{benchmark\_v2\_1.py}，不比较两个版本后择优。原\texttt{benchmark.py}原字节保留为初始化版本与共享校验、统计、校准库；新入口直接复用这些函数。原快照、精确差异、修订时间和哈希保存于\texttt{research/}，新评估协议同时记录原研究manifest及新评估器版本。研究工具、数值求解器、模型轨迹、全部种子与指标均未修改。

\begin{table}[tbp]
  \centering
  \caption{默认数据协议。表中为设计参数，不是性能结果。}
  \label{tab:dataset-protocol}
  \begin{tabular}{p{0.16\linewidth}p{0.24\linewidth}p{0.29\linewidth}r}
    \toprule
    数据用途 & 场景与抽样方式 & 生成种子 & 案例数\\
    \midrule
    发现集 & 四个核心场景，各16例 & 46191 & 64\\
    开发验证 & 四个核心场景，各24例 & 57383 & 96\\
    主测试 & 每个种子下四场景各30例 & 814711、826013、837901 & 360\\
    压力测试 & 五个压力场景，各40例 & 849391 & 200\\
    独立校准 & 四核心场景的独立均匀混合 & 862193 & 240\\
    覆盖评价 & 同分布的独立均匀混合 & 875149 & 320\\
    \bottomrule
  \end{tabular}
\end{table}

四个核心场景分别为静态正常、静态异常、恒速正常和恒速异常。五个压力场景为狭窄几何、加速目标、稀疏观测、标定失配和高污染。主测试采用三个数据种子，以增加独立案例覆盖；这不是三次独立的语言模型研究或规划种子重复，不能据此估计研究控制器的随机性。压力测试有意改变生成机制，用于观察失效边界，不能与同分布主测试合并成一个不加说明的成功率。

冻结记录包含所选配置及其开发验证证据。测试端随后自动产生的比较方法没有反馈给第一轮研究控制器，协议禁止根据其结果修改该冻结方案。早期规则原型的结果已成为第一轮设计的已知背景；本报告后半部分的第二轮又把第一轮结果作为先验，每次后续研究均使用不同的正式测试种子。这里的盲测指工具访问与执行顺序的隔离，并非对拥有完整项目文件的操作者进行密码学保密。可审查性来自冻结文件、协议、输入与预测的保存以及实际事件顺序。

所有数据均由公开模拟器生成。标定失配和高污染只代表预设的合成压力，不是从真实无线设备提取的统计规律。本次没有外部真实数据集或原工业验收数据，任何分场景改善都仅在给定生成机制与样本规模下成立。

\subsection{基线、结构消融与开发前驱}

点估计比较包含四个保留基线：静态WLS、固定参数的静态稳健模型、固定参数的恒速稳健模型，以及第一版的启发式BIC门控。它们由原定位实现包装得到，增加公共诊断不改变已有位置预测。主要比较对象为语言模型最终冻结的方法，因而总体效果同时包含残差、尺度、拟合与模型选择的共同作用。

为区分组件作用，评估脚本依据冻结配置机械生成适用的消融：将对数残差改回距离残差、将常数尺度改回报告尺度、关闭截尾、关闭子集初始化、关闭对数中心修正、将交叉验证门控改为固定恒速，以及将非静态方法改为固定静态。只有所选方案确实包含相应组件时才建立该比较；全部参数变化在协议中保存，且不在测试后重新调优其余参数。

这些比较回答的是“在当前组合附近去掉一个组件会怎样”，不是各个方法族经过公平独立调参后的最优性能比较。例如，更换残差但保持同一稳健尺度，会同时改变该数值尺度的实际含义；若消融表现较差，不能据此断言整个距离残差方法族较差。固定恒速或固定静态也改变计算量，其差异应与工作量共同解释。

脚本还保留登记表中最早两个新增方法作为开发前驱；若配置与最终方法相同则不重复加入。它们不是从最终测试中挑选的最佳中间结果，也不是额外的独立研究运行。前驱比较用于呈现方法修订轨迹，不能替代真实工具日志或单独证明语言模型规划优于其他规划策略。

\subsection{误差统计与配对不确定性}

对每个案例，以当前时刻的位置欧氏误差
\begin{equation}
 e_{m,j}=\|\widehat x_{m,j}(0)-x_j(0)\|_2
\end{equation}
作为基本结果。保存的指标包括均值、RMSE、中位数、P90、P95、最大误差、5米内比例和30米内比例；P90与P95使用NumPy默认的线性分位数插值。文中的CEP90特指这组误差半径的经验90\%分位数。它不等同于原题未定义的“检测精度”，也不能代替静态场景非劣判定。

所有方法使用相同案例。令 $a$ 为冻结方法，$b$ 为比较方法，在场景 $s$ 上定义配对改善
\begin{equation}
 d_{s,j}=e_{b,s,j}-e_{a,s,j},\qquad
 \widehat\Delta=\frac1S\sum_{s=1}^{S}\frac1{n_s}
 \sum_{j=1}^{n_s}d_{s,j}.
 \label{eq:paired-macro-improvement}
\end{equation}
正值表示冻结方法的平均误差较低。\texttt{paired\_ci}在每个核心场景内对差值有放回重采样，再对场景均值等权平均；主测试三个数据种子在各场景内合并，不把三个种子均值当作仅有的三个观测。默认使用种子612303进行3000次重采样，取2.5\%与97.5\%分位数形成名义95\%百分位区间。

这个区间只反映固定方法在当前场景分层与案例抽样下的平均误差差值不确定性。它不是CEP90差异的置信区间，不包含语言模型研究重复、候选选择、实际无线环境迁移或所有计算环境变化。比较多条消融与前驱时，当前实现没有多重比较校正，也没有预注册的单一确证性检验。因此，区间用于探索性的效应量解释；不能把若干区间中偶然较有利的一条称为经过多重检验控制的确认性发现。总体改善仍需结合各静态、移动场景的退化检查。

\subsection{独立校准、覆盖与拒绝曲线}

预测区域仅对已冻结的所选方法进行校准。校准集与覆盖评价集分别独立选择四个核心场景中的一个，再生成案例；这与主测试固定配额的分层设计不同，目的是匹配案例级交换性的条件。标定尺度函数在校准之前已经由定位实现确定，不利用校准真值反复修改。

校准分数与圆盘沿用式\eqref{eq:conformal-score}，考察 $\alpha=0.10$ 和 $\alpha=0.05$ 两个预设水平。实现使用第 $\lceil(m+1)(1-\alpha)\rceil$ 个顺序统计量，而不是线性插值分位数；样本量不足时半径量化因子为正无穷。输出包括总体经验覆盖、平均与中位半径以及分场景覆盖。相应有限样本边际保证依赖交换性与冻结前提\cite{lei2018distributionfree}；一次有限覆盖测试低于名义水平，不会自动否定该边际定理，反过来经验覆盖良好也不能证明前提在真实部署成立。

同一校准因子还应用于五个压力场景，以展示分布变化后的失准，不对该群体宣称分布无关覆盖。拒绝规则的半径门槛预设为2、5、10、15、20、30、50米，分别报告接受率、接受案例数、接受样本平均误差与接受样本覆盖；无接受样本时后两项为空。阈值曲线用于展示覆盖范围与误差的取舍，不作为测试后选择部署阈值的验证依据。尤其不能仅挑一个误差较小但拒绝了大多数案例的点作为总体定位成绩。

当前校准层没有同时为所有基线建立预测圆盘。因此，点估计方法间的不确定性比较由共同案例上的配对区间支持，区域覆盖只评价冻结方法自身；不能据此宣称它的区域比其他方法更窄或更可靠。普通边际覆盖同样不保证拒绝后被接受子集的条件覆盖。

\subsection{计算成本与证据保存}

新求解器的\texttt{work\_units}累计优化器内部实际残差函数调用，包含多个起点、子集求解、截尾轮次、交叉验证折和最终全量重拟合；旧基线保留原求解器记录的函数评价次数。它们都是工作量代理，而非浮点运算总量：一次调用可能处理不同数量的观测，优化器外的评分也不在该计数内。解析Jacobian评价次数与求解次数保存在诊断中，实际耗时另以墙钟时间记录。仅比较方法数量或工具调用次数不足以证明计算预算公平。

基准总耗时包括各方法主测试、压力评价、校准与结果写入，不包含此前的全部语言模型规划、审阅和开发实验成本。科研系统的端到端成本因此需要另读运行时证据，不能把单次定位耗时当作完整研究开销。跨机器耗时差异还受数值库、硬件和系统负载影响。

协议保存冻结文件的字节摘要、定位与工具代码及基准脚本的字节摘要、软件环境、完整输入和逐案例预测。数值求解若发生无法处理的错误，当前评估会中止而非悄悄删除案例；只有全部指定案例完成后，相关汇总才构成完整评估。审计关注的是配置、数据、输出和结论能否相互对应，而不是要求时间戳或不同硬件上的耗时逐字节一致。
